%% =====================================================================
\section{Introduction}
\label{sec:intro}

The \Cu{N} Rubik's Cube is a standard testbed for planning on graphs that are far too large to
enumerate. Its worst-case solution length, God's number, is 20 face turns for $N=3$~\cite{rokicki2014diameter}
(26 quarter turns~\cite{rokicki2014qtm}) and grows as $\Theta(N^2/\log N)$~\cite{demaine2011algorithms};
finding a shortest solution of a given state is NP-complete~\cite{demaine2018npcomplete}. Learned solvers
have recently become strong on small cubes. DeepCubeA~\cite{agostinelli2019deepcubea} and
EfficientCube~\cite{takano2023selfsupervision} solve the \Cu{3}, and CayleyPy~\cite{chervov2025cayleypy}
reaches near-optimal \Cu{3} solutions and the first learned \Cu{4} and \Cu{5} solvers. All of them train one
network per size, and none reports a cube beyond $N=5$. At the other end of the scale, the strongest public
large-cube solutions come from a programming competition, Kaggle's Santa 2023~\cite{kaggle2023santa}, where
the team KMCoders solved, among others, a \Cu{33} in 4,719 quarter turns~\cite{kmcoders2024santa}.

We ask what it takes to cover the whole range of sizes, for both the cube and its flat relative, the \Sq{N}
Rubik's Square of Demaine et al.~\cite{demaine2011algorithms}; which approach gives the shortest solutions at
each size; whether any method can beat the published competition solutions; and how far all of them are
from optimal. To answer on common ground, we run eight solvers and three external baselines on one open engine,
on the same seeded random Squares and the same Santa 2023 cube puzzles; every number in this paper is generated
by script from the result files.

\paragraph{Contributions.}
\begin{itemize}
\item \textbf{\Mseven, a symmetry portfolio} (\S\ref{sec:m7}), our main method. Around KMCoders' program it is the best
  method at every Santa 2023 cube size from $N=4$ to 7, ahead of their published best-of-many solutions on average,
  and second to them at $N=8$ and 9 (Table~\ref{tab:maincube}); we show why its gain shrinks with $N$ (\S\ref{sec:m7scale}).
  Around our strongest Square solvers it is the best method at \MsevenRedSquare{} of the \SizesSquare{} Square sizes
  (Table~\ref{tab:mainsq}).
\item \textbf{\FiveC, net-gain scheduling of bulk commutators} (\S\ref{sec:5c}): the strongest single solver on large
  Squares, and \FiveCGainTwenty\% shorter than Demaine's algorithm on the random \Cu{20}.
\item \textbf{\FiveA, one network for every size} (\S\ref{sec:5a}): a cluster-token transformer that solves Square
  sizes it never saw in training.
\item \textbf{How far from optimal} (\S\ref{sec:bound}): a counting lower bound for every \Cu{N}, calibrated on the
  first exact God's numbers we know of for the \Sq{4} and \Sq{5}, and measured against every method.
\item \textbf{An open engine and benchmark} (\S\ref{sec:engine}) for both puzzles at every $N$, with a harness that
  re-applies and verifies every solution and an exact adapter for the Santa 2023 puzzles.
\end{itemize}
We also reproduce five known methods: \Mone~\cite{korf1997optimal}, \Mtwo~\cite{kociemba2phase},
\Mthree~\cite{demaine2011algorithms}, \Mfour~\cite{chervov2025cayleypy} and \Msix, the KMCoders pipeline rebuilt from its
write-up~\cite{kmcoders2024santa}. We run three external baselines unchanged (\S\ref{sec:baselines}): \Bone, \Btwo and
\Bthree~\cite{walton2017nxnxn}. Table~\ref{tab:names} lists every name.
